% Pandoc LaTeX template for lecture summaries (xelatex).
% Usage: see .claude/commands/lecture.md, step 7.
\documentclass[11pt,a4paper]{article}

% ---------- page & fonts ----------
\usepackage[a4paper,top=2.3cm,bottom=2.2cm,left=2.2cm,right=2.2cm,headsep=0.8cm]{geometry}
\usepackage{fontspec}
\setmainfont{PT Serif}
\setsansfont{PT Sans}
\setmonofont{DejaVu Sans Mono}[Scale=0.88]
\usepackage{polyglossia}
\setdefaultlanguage{russian}
\setotherlanguage{english}
\usepackage{microtype}
\setlength{\parindent}{1.2em}
\setlength{\parskip}{0.35ex plus 0.2ex}
\linespread{1.06}
\frenchspacing
\emergencystretch=3em
\clubpenalty=10000
\widowpenalty=10000
\displaywidowpenalty=10000
% inline formulas (y = 0) are not broken at relations or operators
\relpenalty=10000
\binoppenalty=10000

% ---------- colors ----------
\usepackage[table]{xcolor}
\definecolor{accent}{HTML}{1F4E79}
\definecolor{accent2}{HTML}{C0504D}
\definecolor{accent3}{HTML}{2E7D5B}
\definecolor{muted}{HTML}{6B7280}
\definecolor{paper}{HTML}{F4F7FB}
\definecolor{accent4}{HTML}{9A5B12} % lab primer: "pitfall" boxes

% ---------- packages ----------
\usepackage{amsmath,amssymb}
% amssymb has no blackboard digits: \mathbb{1} (indicator) is drawn as a double-struck one
\let\amsmathbb\mathbb
\renewcommand{\mathbb}[1]{\ifnum\strcmp{#1}{1}=0 \mathrm{1\mkern-4mu l}\else\amsmathbb{#1}\fi}
\usepackage{tikz}
\usepackage[most]{tcolorbox}
\usepackage{titlesec}
\usepackage{fancyhdr}
\usepackage{enumitem}
\usepackage{longtable,booktabs,array,calc}
\usepackage{needspace}
\usepackage{graphicx}
\usepackage{float}
\floatplacement{figure}{H}
\usepackage[font={small,sf},labelformat=empty,justification=centering,skip=4pt]{caption}
\DeclareCaptionFont{muted}{\color{muted}}
\captionsetup{textfont=muted}
\usepackage{etoolbox}
\usepackage{newunicodechar}
% second argument of \texorpdfstring keeps the character itself in PDF bookmarks
\newunicodechar{→}{\texorpdfstring{\ensuremath{\rightarrow}}{→}}
\newunicodechar{←}{\texorpdfstring{\ensuremath{\leftarrow}}{←}}
\newunicodechar{×}{\texorpdfstring{\ensuremath{\times}}{×}}
\newunicodechar{Σ}{\texorpdfstring{\ensuremath{\Sigma}}{Σ}}
\newunicodechar{α}{\texorpdfstring{\ensuremath{\alpha}}{α}}
\newunicodechar{β}{\texorpdfstring{\ensuremath{\beta}}{β}}
\newunicodechar{λ}{\texorpdfstring{\ensuremath{\lambda}}{λ}}
\newunicodechar{ᵀ}{\texorpdfstring{\ensuremath{^{\mathsf{T}}}}{ᵀ}}
\usepackage{hyperref}
\hypersetup{colorlinks=true,linkcolor=accent,urlcolor=accent,linktoc=all,
  bookmarksopen=true,bookmarksopenlevel=1,bookmarksnumbered=true,pdfpagemode=UseOutlines,
  pdftitle={$title$},pdfauthor={$lecturer$},pdfsubject={$course$}}
\usepackage{bookmark}

% ---------- pandoc glue ----------
\providecommand{\tightlist}{\setlength{\itemsep}{0pt}\setlength{\parskip}{0pt}}
% images: never wider than the text or taller than 45% of the page, with a thin frame
\makeatletter
\newsavebox\pandoc@box
% (both limits are checked in turn: a tall portrait board scan, scaled to the text width,
% would otherwise fill a whole page)
\newcommand{\pandocbounded}[1]{\sbox\pandoc@box{#1}%
  \ifdim\wd\pandoc@box>0.94\linewidth\sbox\pandoc@box{\resizebox{0.94\linewidth}{!}{\usebox\pandoc@box}}\fi
  \ifdim\ht\pandoc@box>0.45\textheight\sbox\pandoc@box{\resizebox{!}{0.45\textheight}{\usebox\pandoc@box}}\fi
  \usebox\pandoc@box}
\makeatother
\setkeys{Gin}{keepaspectratio}
\newcounter{none} % pandoc: uncaptioned longtables use \LTcaptype{none}
\setlength{\LTpre}{0.8em}\setlength{\LTpost}{0.8em}
\arrayrulecolor{accent!70}
\setlength{\heavyrulewidth}{0.9pt}\setlength{\lightrulewidth}{0.4pt}
\renewcommand{\arraystretch}{1.18}
\AtBeginEnvironment{longtable}{\small}
\newcommand{\tablehead}[1]{{\sffamily\bfseries\color{accent}#1}}
% code blocks with a language (pandoc highlighting): long lines wrap, the block sits in a frame
% (codeframe); "Плохо:"/"Хорошо:" pairs recolor the frame via codevariant (see conspect.lua)
$if(highlighting-macros)$
$highlighting-macros$
\usepackage{fvextra}
\RecustomVerbatimEnvironment{Highlighting}{Verbatim}{commandchars=\\\{\},breaklines,breakanywhere,fontsize=\small}
\renewenvironment{Shaded}{\begin{codeframe}}{\end{codeframe}}
$endif$

% ---------- headings ----------
\setcounter{secnumdepth}{1}
\titleformat{\section}
  {\sffamily\Large\bfseries\color{accent}}
  {\tikz[baseline=(n.base)]\node[fill=accent,text=white,rounded corners=3pt,inner xsep=6pt,inner ysep=3pt,minimum width=1.6em](n){\thesection};}
  {0.6em}{}
  [\vspace{-0.6ex}{\color{accent!25}\titlerule[1pt]}]
% running head for unnumbered sections too; \phantomsection gives the following \addcontentsline
% (pandoc: \section*{..}\addcontentsline) its own anchor, otherwise TOC links/bookmarks of the
% unnumbered sections point to the previous numbered heading
\newcommand{\markunnumbered}[1]{\phantomsection\markboth{#1}{}#1}
\titleformat{name=\section,numberless}
  {\sffamily\Large\bfseries\color{accent}}{}{0pt}{\markunnumbered}
  [\vspace{-0.6ex}{\color{accent!25}\titlerule[1pt]}]
\titlespacing*{\section}{0pt}{3.2ex plus 1ex minus .5ex}{1.4ex}
% ---------- table of contents (clickable; PDF outline comes from hyperref/bookmark) ----------
\usepackage{titletoc}
\setcounter{tocdepth}{2}
% video timecode + slides of a section are shown on its TOC line: \secmeta stores them in the .aux
% (\tocmetadef{<section no>}{<timecode>}{<slides>}), the section entry reads them back.
% section entry: number | title (wraps in its own column) | video+slides (right column) | page
% (13.6em: "слайды 20 · доска 22–29" fits on one line)
\newlength{\tocmetawidth}\setlength{\tocmetawidth}{13.6em}
\newcommand{\tocsectitle}[1]{\parbox[t]{\dimexpr\linewidth-\tocmetawidth-6.2em\relax}{\raggedright #1}}
% density of the TOC; \tocpage makes it more compact if the TOC does not fit on one page
\newcommand{\tocsecskip}{0.2em}
\newcommand{\tocsubfont}{\fontsize{10}{11}\selectfont}
\titlecontents{section}[2.4em]{\addvspace{\tocsecskip}\sffamily\bfseries\color{accent}}
  {\contentslabel[\thecontentslabel]{2.4em}\tocsectitle}{\hspace*{-2.4em}\tocsectitle}
  {\hfill\parbox[t]{\tocmetawidth}{\raggedleft\leavevmode\tocmeta{\thecontentslabel}}\makebox[2em][r]{\contentspage}}
  [\addvspace{1.5pt}]
% last-resort subsection format for an overfull TOC (see \tocpage): subsections run in one
% paragraph per section, "title page · title page · …", still clickable. After \begin{document}
% \titlecontents only writes a format change into the .toc (applied on the next pass, which makes
% the one-page check oscillate), so both formats are defined here and \tocrunin switches them.
\titlecontents*{subsection}[2.4em]{\tocsubfont\rightskip=0pt plus 4em}
  {}{}{\nobreak\enspace\textcolor{muted}{\thecontentspage}}[\enspace\textperiodcentered\enspace][\tocrunend]
\makeatletter
\let\tocrunin@l\l@subsection
\expandafter\let\expandafter\tocrunin@b\csname ttlb@subsection\endcsname
\expandafter\let\expandafter\tocrunin@m\csname ttlm@subsection\endcsname
\expandafter\let\expandafter\tocrunin@e\csname ttle@subsection\endcsname
\newcommand{\tocrunin}{\let\l@subsection\tocrunin@l
  \expandafter\let\csname ttlb@subsection\endcsname\tocrunin@b
  \expandafter\let\csname ttlm@subsection\endcsname\tocrunin@m
  \expandafter\let\csname ttle@subsection\endcsname\tocrunin@e}
\makeatother
\newcommand{\tocrunend}{\par\addvspace{2pt}}% \par is not allowed in \titlecontents* args
\titlecontents{subsection}[2.4em]{\tocsubfont}
  {}{}{\titlerule*[0.7pc]{.}\contentspage}
\makeatletter
\newcommand{\tocmetadef}[3]{\expandafter\gdef\csname tocmeta@#1\endcsname{#2\enspace #3}}
\newcommand{\tocmeta}[1]{\ifcsname tocmeta@#1\endcsname
  {\mdseries\footnotesize\color{muted}\playicon\ \csname tocmeta@#1\endcsname}\fi}
\makeatother
\newcommand{\tocvideo}[2]{} % obsolete .toc entries from older builds
% The TOC must stay on one page: typeset the .toc (from the previous pass) into a box and,
% while it is taller than the page, switch to denser settings. The box comes out ~30pt taller
% than the typeset TOC, and \enlargethispage gives 2 more lines; the threshold is \textheight
% (\textheight+1 line was calibrated on lectures/r2, but a 729.99pt box still overflowed by one
% line on lectures/r_seminar_1, where long section meta wraps to several lines).
\newsavebox{\tocmeasurebox}
\makeatletter
\newcommand{\tocmeasure}{\IfFileExists{\jobname.toc}{\sbox{\tocmeasurebox}{\vbox{\hsize\textwidth
  \setlength{\parskip}{0pt}\makeatletter\@input{\jobname.toc}}}}{\sbox{\tocmeasurebox}{}}}
\makeatother
\newcommand{\toctoobig}{\dimexpr\ht\tocmeasurebox+\dp\tocmeasurebox\relax>\textheight}
\newcommand{\tocpage}{\clearpage\pdfbookmark[1]{Оглавление}{toc}%
  \renewcommand{\contentsname}{Оглавление}{\setlength{\parskip}{0pt}\linespread{1}\selectfont
  \tocmeasure
  \ifdim\toctoobig\renewcommand{\tocsecskip}{0.05em}\renewcommand{\tocsubfont}{\fontsize{9.5}{10.3}\selectfont}\tocmeasure\fi
  \ifdim\toctoobig\renewcommand{\tocsecskip}{0pt}\renewcommand{\tocsubfont}{\fontsize{9}{9.6}\selectfont}\tocmeasure\fi
  \ifdim\toctoobig\renewcommand{\tocsubfont}{\fontsize{8.5}{9}\selectfont}\tocmeasure\fi
  % last resort (many subsections, e.g. lectures/p3 with 68): run-in subsections
  \ifdim\toctoobig\renewcommand{\tocsubfont}{\fontsize{9.5}{11.2}\selectfont}\tocrunin\fi
  \enlargethispage{2\baselineskip}\tableofcontents}\clearpage}

\titleformat{\subsection}{\sffamily\large\bfseries\color{accent!85!black}}{}{0pt}{}
\titlespacing*{\subsection}{0pt}{2.2ex plus .8ex minus .3ex}{0.8ex}
\pretocmd{\section}{\Needspace{9\baselineskip}}{}{}
\pretocmd{\subsection}{\Needspace{5\baselineskip}}{}{}

% video timecode + slides under a section title
\newcommand{\playicon}{\tikz[baseline=-0.1ex]\fill[accent] (0,-0.05ex) -- (0,1.55ex) -- (1.25ex,0.75ex) -- cycle;}
\makeatletter
\newcommand{\secmeta}[2]{\par\noindent
  \tikz[baseline=(t.base)]\node[fill=accent!9,rounded corners=2.5pt,inner xsep=5pt,inner ysep=2.5pt](t)
    {\sffamily\small\playicon\ \textcolor{accent}{\bfseries #1}};%
  \hspace{0.6em}{\sffamily\small\color{muted}#2}\par\medskip
  \protected@write\@auxout{}{\string\tocmetadef{\thesection}{#1}{#2}}}
\makeatother

% inline labels
\newcommand{\compnote}{{\sffamily\bfseries\color{accent2}Прим.\,составителя.}\ }
\newcommand{\lectmark}[1]{{\sffamily\bfseries\color{accent2}#1}}
\newcommand{\goodmark}[1]{{\sffamily\bfseries\color{accent3}#1}}

% ---------- lists ----------
\setlist{itemsep=0.25ex,topsep=0.6ex,parsep=0pt}
\setlist[itemize,1]{label={\color{accent}\small$$\blacktriangleright$$},leftmargin=1.5em}
\setlist[itemize,2]{label={\color{accent}\textendash},leftmargin=1.3em}
\setlist[enumerate,1]{label={\sffamily\bfseries\color{accent}\arabic*.},leftmargin=1.8em}

% ---------- boxes ----------
\tcbset{conspect/.style={enhanced,breakable,sharp corners=south,arc=4pt,
  boxrule=0pt,leftrule=3pt,left=10pt,right=10pt,top=7pt,bottom=7pt,
  fonttitle=\sffamily\bfseries,coltitle=white}}
\newtcolorbox{abstractbox}[1]{conspect,colback=paper,colframe=accent,leftrule=3pt,
  title={#1},attach title to upper,coltitle=accent,fonttitle=\sffamily\bfseries,after title={\par\smallskip},
  fontupper={\setlength{\parskip}{0.8ex}}}
\newtcolorbox{keypointsbox}[1]{conspect,colback=accent!5,colframe=accent}
\newtcolorbox{emphasisbox}[1]{conspect,colback=accent2!5,colframe=accent2}
\newtcolorbox{watchbox}[1]{conspect,colback=accent3!6,colframe=accent3}
\newtcolorbox{qabox}[1]{enhanced,breakable,colback=muted!6,colframe=muted!55,boxrule=0pt,leftrule=2pt,
  arc=0pt,left=9pt,right=9pt,top=5pt,bottom=6pt,
  fontupper={\small\setlength{\parskip}{0.6ex}\setlength{\parindent}{0pt}},
  title={\playicon\hspace{0.2em}\textsf{\bfseries #1}},attach title to upper,coltitle=muted,after title={\par\smallskip}}
% breakable boxes with the title attached to the text: don't leave the title alone at the page bottom
\BeforeBeginEnvironment{qabox}{\Needspace{5\baselineskip}}
\BeforeBeginEnvironment{labnotebox}{\Needspace{4\baselineskip}}
\BeforeBeginEnvironment{labwarnbox}{\Needspace{4\baselineskip}}
\BeforeBeginEnvironment{labcautionbox}{\Needspace{4\baselineskip}}
\BeforeBeginEnvironment{labaskbox}{\Needspace{4\baselineskip}}
% lab primer: GitHub alerts (> [!NOTE] …) -> boxes, see conspect.lua
\tcbset{labbox/.style={conspect,title={#1},attach title to upper,after title={\par\smallskip},
  fontupper={\setlength{\parindent}{0pt}\setlength{\parskip}{0.6ex}}}}
\newtcolorbox{labnotebox}[1]{labbox={#1},colback=accent!5,colframe=accent,coltitle=accent}
\newtcolorbox{labwarnbox}[1]{labbox={#1},colback=accent4!6,colframe=accent4,coltitle=accent4}
\newtcolorbox{labcautionbox}[1]{labbox={#1},colback=accent2!5,colframe=accent2,coltitle=accent2}
\newtcolorbox{labaskbox}[1]{labbox={#1},colback=muted!6,colframe=muted!55,coltitle=muted,leftrule=2pt,
  fontupper={\small\setlength{\parindent}{0pt}\setlength{\parskip}{0.6ex}}}
\tcbset{codevariant/.style={}}
\newtcolorbox{codeframe}{enhanced,breakable,colback=paper,colframe=accent!50,boxrule=0pt,leftrule=2.5pt,
  arc=0pt,left=6pt,right=6pt,top=3pt,bottom=3pt,before skip=0.8ex,after skip=0.8ex,codevariant}
% caption under a Mermaid scheme (*Схема: …*)
\newenvironment{schemecap}{\par\begingroup\small\sffamily\color{muted}\centering}{\par\endgroup\medskip}
\newtcolorbox{planbox}[1]{conspect,colback=white,colframe=accent!60,boxrule=0.4pt,leftrule=0.4pt,
  left=4pt,right=4pt,top=2pt,bottom=2pt,arc=3pt,sharp corners=false}

% ---------- footnotes ----------
\renewcommand{\footnoterule}{\vspace*{-3pt}{\color{accent!40}\hrule width 0.3\columnwidth height 0.5pt}\vspace*{2.6pt}}
\makeatletter
% the mark sits in a 1.2em column; two-digit marks (10, 11, ...) get a space instead of touching the text
\renewcommand{\@makefntext}[1]{\noindent\hangindent=1.2em\sbox0{\color{accent2}\sffamily\bfseries\@thefnmark\hspace{0.3em}}%
  \ifdim\wd0<1.2em\hbox to 1.2em{\unhbox0\hfil}\else\box0\fi#1}
\renewcommand{\thefootnote}{\arabic{footnote}}
\makeatother

% ---------- header / footer ----------
\pagestyle{fancy}
\fancyhf{}
\renewcommand{\sectionmark}[1]{\markboth{#1}{}}
\fancyhead[L]{\sffamily\footnotesize\color{muted}$if(lab)$$lab$$else$$lecture$$endif$}
\fancyhead[R]{\sffamily\footnotesize\color{muted}\parbox[b]{0.8\textwidth}{\raggedleft\nouppercase{\leftmark}}}
\fancyfoot[C]{\sffamily\footnotesize\color{muted}\thepage}
\renewcommand{\headrulewidth}{0.4pt}
\renewcommand{\headrule}{{\color{accent!20}\hrule height\headrulewidth width\headwidth}}
\fancypagestyle{first}{\fancyhf{}\renewcommand{\headrulewidth}{0pt}
  \fancyfoot[C]{\sffamily\footnotesize\color{muted}\thepage}}

\begin{document}
\thispagestyle{first}

% ---------- hero ----------
% the band grows with the title (a 3-line title must not push the subtitle out of it)
\newsavebox{\herobox}\newlength{\heroband}
\sbox{\herobox}{\parbox[t]{16.6cm}{\raggedright
    {\sffamily\color{white!80!accent}\large $course$\par}\vspace{0.9em}
    {\sffamily\bfseries\color{white}\fontsize{27}{31}\selectfont\hyphenpenalty=10000 $title$\par}\vspace{0.8em}
    {\sffamily\color{white!85!accent}\large $if(lab)$$lab$ · вопросы к защите$else$$lecture$ · конспект$if(lecturer)$ · $lecturer$$endif$$endif$\par}}}
\setlength{\heroband}{\dimexpr\ht\herobox+\dp\herobox+3.1cm\relax}
\ifdim\heroband<6.4cm\setlength{\heroband}{6.4cm}\fi % 6.4cm band, 4.9cm skip: the original layout
\begin{tikzpicture}[remember picture,overlay]
  \fill[accent] (current page.north west) rectangle ([yshift=-\heroband]current page.north east);
  \fill[accent2] ([yshift=-\heroband]current page.north west) rectangle ([yshift=-\heroband-0.15cm]current page.north east);
  \node[anchor=north west,inner sep=0pt]
    at ([xshift=2.2cm,yshift=-1.55cm]current page.north west) {\usebox{\herobox}};
\end{tikzpicture}
\vspace*{\dimexpr\heroband-1.5cm\relax}

$if(meta-line)$
{\noindent\sffamily\small\color{muted}$meta-line$\par}
\bigskip
$endif$

$body$

\end{document}
